% ============ 第4章 §4.8 Kakeya集的几何维数(原书 pp.367-371) ============
\section{Kakeya 集的几何维数}\label{sec:4.8}
\index{Kakeya 集}
\index{Besicovitch 集}
\index{Hausdorff 外测度}
\index{Hausdorff 上维数}
\index{Minkowski 外测度}
\index{Minkowski 上维数}
\index{tube!tube 的维数}


本节作为前一节定理~\ref{thm:4.7.1} 的一个应用,讨论 Kakeya 集的几何维数的估计.首先定义 Kakeya 集.

\begin{definition}\label{def:4.8.1}
$\mathbb{R}^{n}$($n > 1$)中 $n$ 维 Lebesgue 测度为 0 的紧致子集,若其内部包含一切方向的单位线段,则称之为 Kakeya 集或 Besicovitch 集.
\end{definition}

关于这种令人着迷的集合的构造方法,请参考其他书籍.这里,我们只考虑该集合的几何维数的估计.

来定义几何维数.

\begin{definition}\label{def:4.8.2}
对 $E \subset \mathbb{R}^{n}$,其 Hausdorff 外测度定义为
\begin{equation*}
\mathcal{H}_{s}(E) := \sup_{\delta > 0}\inf\left\{\sum_{j}t_{j}^{s}:\ E \subset \bigcup_{j}B(x_{j},t_{j}),\ t_{j} < \delta\right\} \quad (s > 0).
\end{equation*}
\end{definition}

对有界的无限集合 $E \subset \mathbb{R}^{n}$,下式成立:
\[
\mathcal{H}_{n+\epsilon}(E) = 0 \quad (\epsilon > 0), \qquad \mathcal{H}_{0}(E) = \infty.
\]

\begin{definition}\label{def:4.8.3}
对 $E \subset \mathbb{R}^{n}$,其上 Hausdorff 维数定义为
\begin{equation*}
\operatorname{H-dim}(E) := \inf\{s:\ \mathcal{H}_{s}(E) = 0\}.
\end{equation*}
\end{definition}

\begin{definition}\label{def:4.8.4}
对 $E \subset \mathbb{R}^{n}$,其 Minkowski 外测度定义为
\begin{equation*}
\mathcal{M}_{s}(E) := \limsup_{\delta \to 0}\inf\left\{\sum_{j}\delta^{s}:\ E \subset \bigcup_{j}B(x_{j},\delta)\right\} \quad (s > 0).
\end{equation*}
\end{definition}

由定义得
\begin{equation}\tag{4.8.1}
\mathcal{M}_{s}(E) \geq \mathcal{H}_{s}(E).
\label{eq:4.8.1}\end{equation}

\begin{definition}\label{def:4.8.5}
对 $E \subset \mathbb{R}^{n}$,其上 Minkowski 维数定义为
\begin{equation*}
\operatorname{M-dim}(E) := \inf\{s:\ \mathcal{M}_{s}(E) = 0\}.
\end{equation*}
\end{definition}

对 $1 \leq p \leq n$,把"$\mathbb{R}^{n}$ 上所有 Kakeya 集具有至少 $p$ 的 Hausdorff 维数"这一主张记为 $K_{H}(p,n)$,把"$\mathbb{R}^{n}$ 上所有 Kakeya 集具有至少 $p$ 的 Minkowski 维数"这一主张记为 $K_{M}(p,n)$.此时,如下猜想成立.

\textbf{猜想问题.} $K_{H}(n,n)$ 以及 $K_{M}(n,n)$ 成立.

由\eqref{eq:4.8.1}可知 $K_{H}(p,n) \Rightarrow K_{M}(p,n)$ 的关系.

下面说明上一节讨论的小 Nikodym 极大函数 $K_{\delta}f$ 的猜想问题与上面的猜想问题之间的联系.

对 $1 < p \leq n$,把"上一节的\eqref{eq:4.7.5}在 $p$ 时成立"这一主张记为 $K_{N}(p,n)$.重复上一节的主张可知 $K_{N}(n,n)$ 被猜想.特别地,由定理~\ref{thm:4.7.1} 知 $K_{N}((n+2)/2, n)$ 成立.

\begin{lemma}\label{lem:4.8.6}
设 $1 < p \leq n$,则
\[
K_{N}(p,n) \Rightarrow K_{M}(p,n), \quad K_{N}(p,n) \Rightarrow K_{H}(p,n)
\]
的关系成立.
\end{lemma}

\textbf{$K_{N}(p,n) \Rightarrow K_{M}(p,n)$.} 对 Kakeya 集 $E \subset \mathbb{R}^{n}$,当 $s < p$ 时证明 $\mathcal{M}_{s}(E) > 0$.

设 $0 < \delta \ll 1$ 为小参数,把 $E$ 的 $\delta$ 邻域记为 $\mathcal{N}_{\delta}(E)$.此时,若把 $\mathcal{N}_{\delta}(E)$ 的特征函数记为 $f$,则由简单的几何考察可知
\begin{align*}
\|K_{\delta}f\|_{q} &> c > 0. \\
\intertext{用此式应用假设,则对某个 $\alpha$,}
\delta^{n-p}(\log(1/\delta))^{-\alpha} &\leq C|\mathcal{N}_{\delta}(E)|.
\end{align*}
因此,覆盖 $E$ 的球 $B(x_{j}, \delta)$ 的个数可由下界 $C^{-1}\delta^{\epsilon-p}$($\epsilon > 0$)估计,当 $s < p$ 时
\[
\mathcal{M}_{s}(E) \geq \limsup_{\delta \to 0}C^{-1}\delta^{-(p-s)/2} = \infty.
\]

\textbf{$K_{N}(p,n) \Rightarrow K_{H}(p,n)$.} 这次对 Kakeya 集 $E$,当 $s < p$ 时证明 $\mathcal{H}_{s}(E) > 0$.设 $E \subset \bigcup_{j}B(x_{j},t_{j})$,$t_{j} < 1$,令
\[
E_{k} := \bigcup_{j:\ 2^{-k} \leq t_{j} < 2^{1-k}}B(x_{j},t_{j}) \quad (k = 1, 2, \dots).
\]
由 Kakeya 集的性质,对任意 $\omega \in S^{n-1}$,存在包含于 $E$ 的单位线段 $l_{\omega}$.把一维 Hausdorff 测度记为 $|\cdot|$,则
\[
1 = |l_{\omega} \cap E| \leq \sum_{k}|l_{\omega} \cap E_{k}| = \sum_{k}\frac{|l_{\omega} \cap E_{k}|k^{2}}{k^{2}}.
\]
因此,用 $\sum_{k}\frac{1}{k^{2}} = c$,可选取正整数 $k_{\omega}$ 使得
\begin{align*}
|l_{\omega} \cap E_{k_{\omega}}| &> \frac{1}{ck_{\omega}^{2}}. \\
\intertext{同理可选取正整数 $\kappa$,使得}
\left|\left\{\omega \in S^{n-1}:\ |l_{\omega} \cap E_{\kappa}| > \frac{1}{c\kappa^{2}}\right\}\right| &> |S^{n-1}|\frac{1}{c\kappa^{2}}.
\end{align*}
其中左端外侧的 $|\cdot|$ 与右端的 $|\cdot|$ 是 $S^{n-1}$ 的 $n-1$ 维面积测度.

把 $E_{\kappa}$ 的特征函数记为 $f$,取 $\delta = 2^{-\kappa}$,由简单的几何考察,由上式得
\begin{align*}
\|K_{\delta}f\|_{q} &> C(1/\kappa^{2})^{1/q+1}. \\
\intertext{应用假设,则对某个 $\alpha$,}
(2^{-\kappa})^{n-p}\kappa^{-\alpha}(1/\kappa^{2})^{p/q+p} &\leq C|E_{\kappa}|.
\end{align*}
因此,$B(x_{j},t_{j})$($2^{-\kappa} \leq t_{j} < 2^{1-\kappa}$)的个数可由下界 $C^{-1}2^{p\kappa}\kappa^{-\alpha}(1/\kappa^{2})^{p/q+p}$ 估计,
\[
\sum_{j}t_{j}^{s} \geq C^{-1}2^{\kappa(p-s)}\kappa^{-\alpha}(1/\kappa^{2})^{p/q+p} > a > 0.
\]
其中 $a$ 是不依赖于 $\kappa$ 的数.因此 $\mathcal{H}_{s}(E) \geq a > 0$.

由该引理与定理~\ref{thm:4.7.1},如下定理成立.

\begin{theorem}\label{thm:4.8.7}
$\mathbb{R}^{n}$ 的所有 Kakeya 集至少具有 $(n+2)/2$ 的上 Hausdorff 维数以及上 Minkowski 维数.
\end{theorem}
